\chapter{Teoria Morse'a: punkty krytyczne, uchwyty i przepływ}
\label{ch:teoria-morsea}
\index{teoria Morse'a}

Funkcja wysokości pozwala oglądać rozmaitość
po jednej poziomicy naraz. Między jej wartościami
krytycznymi nic istotnego się nie zmienia:
Twierdzenie~\ref{thm:pas-bez-krytycznych-15}
identyfikuje taki pas z iloczynem.
Wartość krytyczna jest chwilą, w której
pojawia się nowa część przestrzeni, otwór
albo zamknięcie otworu. Teoria Morse'a
zamienia ten obraz w twierdzenia o uchwytach
i homologiach. Na końcu poznamy trajektorie
gradientowe, które pozwalają obliczać różniczkę
kompleksu Morse'a.

W głównej części rozdziału $M^n$ jest gładką
rozmaitością \emph{zwartą bez brzegu}, a $f:M\to\R$
funkcją Morse'a. Definicję hesjanu
i niezdegenerowania podaliśmy w
Definicji~\ref{def:morse-pochodna}; istnienie
funkcji Morse'a i możliwość rozdzielenia
ich wartości krytycznych dowiedliśmy w
Twierdzeniu~\ref{thm:gestosc-morse}.
Zatem nie będziemy drugi raz udowadniać
tych samych faktów. Oznaczamy
\[
 M^a:=f^{-1}((-\infty,a]),\qquad
 \operatorname{Crit}_k(f)
  :=\{p:\dd f_p=0,\ \lambda(p)=k\},
 \qquad c_k:=\#\operatorname{Crit}_k(f).
\]
Przy wartości regularnej $a$ zbiór $M^a$
jest zwartą rozmaitością z brzegiem
$f^{-1}(a)$, o ile nie jest pusty.

\section{Lemat Morse'a: dokładny opis w pobliżu punktu}
\label{sec:lemat-morsea-17}

Hesjan w punkcie krytycznym mierzy
drugorzędową zmianę $f$. Jego indeks
$\lambda(p)$, wprowadzony przy
transwersalności w
Sekcji~\ref{sec:transwersalnosc-smale},
jest liczbą ujemnych kierunków.
Lemat Morse'a mówi znacznie więcej
niż wzór Taylora z resztą: po gładkiej
zmianie współrzędnych \emph{reszta znika}.

\begin{twierdzenie}[Lemat Morse'a]
\label{thm:lemat-morsea-17}
Jeśli $p$ jest niezdegenerowanym punktem
krytycznym o indeksie $\lambda$, istnieje
mapa $(u,v)$ wokół $p$, w której
$u\in\R^\lambda$, $v\in\R^{n-\lambda}$,
$p=(0,0)$ i
\begin{equation}\label{eq:postac-morsea-17}
 f(u,v)=f(p)-|u|^2+|v|^2.
\end{equation}
Liczba $\lambda$ nie zależy od mapy.
\end{twierdzenie}
\begin{proof}
Wybierzmy dowolne współrzędne $x$ z $x(p)=0$.
Ponieważ $\dd f_0=0$, liniowa zamiana
współrzędnych, twierdzenie Sylvestera
i przeskalowanie osi pozwalają założyć
\[
 \tfrac12\operatorname{Hess}_0f
  =J:=\operatorname{diag}(-I_\lambda,I_{n-\lambda}).
\]
Nie wolno jednak na tym poprzestać:
forma w punkcie $0$ nie opisuje jeszcze
całej funkcji w sąsiedztwie.
Dwukrotne całkowanie po odcinku
$t\mapsto tx$ daje \emph{dokładną} równość
\[
 f(x)-f(0)
  =x^{\mathsf T}H(x)x,\qquad
 H(x)=\int_0^1(1-t)\,
                    \operatorname{Hess}_{tx}f\,\dd t.
\]
Macierz $H(x)$ zależy gładko od $x$,
jest symetryczna i $H(0)=J$.
Rozbijmy $x=(x_-,x_+)$ i zapiszmy
\[
 H(x)=
 \begin{pmatrix} A(x)&B(x)\\
                 B(x)^{\mathsf T}&C(x)\end{pmatrix}.
\]
Blisko zera $A$ jest ujemnie określona,
a dopełnienie Schura
$S=C-B^{\mathsf T}A^{-1}B$ dodatnio
określone. Istotnie, w zerze mamy
$A=-I$, $B=0$, $S=I$, a określoność
jest własnością otwartą.
Uzupełniamy kwadrat, nie pomijając
zależności macierzy od $x$:
\[
 x^{\mathsf T}Hx
 =(x_-+A^{-1}Bx_+)^{\mathsf T}
       A(x_-+A^{-1}Bx_+)
   +x_+^{\mathsf T}Sx_+.
\]
Dodatnio określone macierze $-A$ i $S$
mają gładko zależne od $x$ dodatnie
pierwiastki kwadratowe. Blisko $I$
można je zbudować z szeregu potęgowego
$\sqrt{I+E}
 =I+\tfrac12E-\tfrac18E^2+\cdots$;
na mniejszym otoczeniu $\|E\|\leq q<1$.
Po dowolnej ustalonej liczbie różniczkowań
wyrazy są ograniczone przez wielomian
w numerze wyrazu razy potęgę $q$.
Szeregi pochodnych są więc lokalnie
jednostajnie zbieżne, co uzasadnia gładkość. Kładziemy
\[
 u=(-A)^{1/2}(x_-+A^{-1}Bx_+),
 \qquad v=S^{1/2}x_+.
\]
Otrzymujemy~\eqref{eq:postac-morsea-17}.
W zerze różniczka zmiany
$x\mapsto(u,v)$ jest identycznością,
bo $A(0)=-I$, $B(0)=0$, $S(0)=I$.
Twierdzenie o funkcji odwrotnej
czyni więc $(u,v)$ współrzędnymi
na pewnym otoczeniu $p$.
Przy $\lambda=0$ pomijamy blok $A$ i bierzemy
$v=H(x)^{1/2}x$; przy $\lambda=n$ bierzemy
$u=(-H(x))^{1/2}x$. Nie odwracamy w tych
przypadkach żadnego nieistniejącego bloku.

Na końcu trzeba sprawdzić jednoznaczność
liczby $\lambda$: hesjan w tych
współrzędnych ma macierz
$\operatorname{diag}(-2I_\lambda,2I_{n-\lambda})$,
a prawo bezwładności Sylvestera
zachowuje liczbę ujemnych kwadratów
przy każdej odwracalnej zmianie bazy.
\end{proof}

\begin{przyklad}[Trzy lokalne kształty w wymiarze dwa]
Na płaszczyźnie funkcje $x^2+y^2$,
$-x^2+y^2$ i $-x^2-y^2$ mają
odpowiednio indeksy $0,1,2$.
Pierwsza tworzy nową składową
podpoziomicy. Przy siodle dwie
części poziomicy stykają się
i zmieniają sposób połączenia.
Ostatnia może wypełnić ostatni
brzeg powierzchni. Sam rysunek
poziomic nie wyznacza jeszcze
globalnego wyniku: o sposobie
przyklejenia decyduje wcześniejsza
część rozmaitości.
\end{przyklad}

\begin{figure}[htbp]
\centering
\begin{tikzpicture}[>=Latex,line cap=round,font=\small,scale=.94]
 \begin{scope}[shift={(-4.15,0)}]
  \fill[gray!5] (0,0) circle (1.35);
  \fill[manifoldblue!40] (0,0) circle (.68);
  \draw[manifoldblue!85!black,very thick] (0,0) circle (.68);
  \draw[gray!50] (0,0) circle (1.35);
  \draw[->,manifoldgold!85!black,thick] (.79,.14)--(1.18,.21);
  \draw[->,manifoldgold!85!black,thick] (-.79,-.14)--(-1.18,-.21);
  \fill[manifoldred] (0,0) circle (2.5pt);
  \node at (0,-1.68) {minimum $\lambda=0$};
  \node[manifoldblue!75!black] at (0,-2.10) {pojawia się dysk};
 \end{scope}
 \begin{scope}
  \begin{scope}
   \clip (0,0) circle (1.35);
   \fill[gray!5] (-1.5,-1.5) rectangle (1.5,1.5);
   \fill[manifoldblue!40] (1.5,-1.5)--(1.5,1.5)
     -- plot[domain=1.5:-1.5,samples=50]
        ({sqrt(\x*\x+.15)},{\x}) -- cycle;
   \fill[manifoldblue!40] (-1.5,-1.5)--(-1.5,1.5)
     -- plot[domain=1.5:-1.5,samples=50]
        ({-sqrt(\x*\x+.15)},{\x}) -- cycle;
   \draw[manifoldgold!85!black,dashed,thick]
     (-1.38,-1.38)--(1.38,1.38)
     (-1.38,1.38)--(1.38,-1.38);
   \foreach \sgn in {-1,1}
    \draw[manifoldblue!85!black,very thick,
          domain=-1.4:1.4,samples=50]
      plot ({\sgn*sqrt(\x*\x+.15)},{\x});
  \end{scope}
  \draw[gray!50] (0,0) circle (1.35);
  \fill[manifoldred] (0,0) circle (2.5pt);
  \node at (0,-1.68) {siodło $\lambda=1$};
  \node[manifoldblue!75!black] at (0,-2.10) {dwa obszary};
 \end{scope}
 \begin{scope}[shift={(4.15,0)}]
  \fill[manifoldblue!40] (0,0) circle (1.35);
  \fill[white] (0,0) circle (.68);
  \draw[manifoldred!80!black,very thick] (0,0) circle (.68);
  \draw[gray!50] (0,0) circle (1.35);
  \draw[->,manifoldgold!85!black,thick] (1.18,.21)--(.79,.14);
  \draw[->,manifoldgold!85!black,thick] (-1.18,-.21)--(-.79,-.14);
  \fill[manifoldred] (0,0) circle (2.5pt);
  \node at (0,-1.68) {maksimum $\lambda=2$};
  \node[manifoldblue!75!black] at (0,-2.10) {otwór się zamyka};
 \end{scope}
\end{tikzpicture}
\caption{Lokalne podpoziomice trzech postaci
z lematu Morse'a w wymiarze dwa. Kolor niebieski
oznacza $f\leq a$: przy minimum pokazujemy $a>f(p)$,
a przy siodle i maksimum $a<f(p)$. Złote strzałki
wskazują wzrost $f$: od minimum na zewnątrz,
do maksimum ku środkowi. Przerywane proste
przy siodle to poziom krytyczny.}
\label{fig:poziomice-morse-17}
\end{figure}

\section{Przekroczenie wartości krytycznej: uchwyt}
\label{sec:uchwyt-morse-17}

Przedział bez wartości krytycznych mamy już
opisany w Twierdzeniu~\ref{thm:pas-bez-krytycznych-15}.
Teraz $a<b$ są wartościami regularnymi,
a w $f^{-1}([a,b])$ leży dokładnie jeden
punkt krytyczny $p$ z $a<f(p)=c<b$ i indeksem
$\lambda$. Dla przejrzystości
piszemy $D^0=\{\mathrm{punkt}\}$
i $S^{-1}=\varnothing$.

\begin{definicja}[Uchwyt indeksu $\lambda$]
\label{def:uchwyt-morse-17}
$n$-wymiarowy \emph{uchwyt indeksu $\lambda$}
to $H_\lambda=D^\lambda\times D^{n-\lambda}$.
Jego część przyczepiana to
$S^{\lambda-1}\times D^{n-\lambda}$.
Rdzeń $D^\lambda\times\{0\}$ ma
wymiar $\lambda$, a dysk
$\{0\}\times D^{n-\lambda}$ biegnie
w kierunkach poprzecznych.
Przyklejamy uchwyt do brzegu
dotychczasowej podpoziomicy
przez zanurzenie jego części przyczepianej.
\end{definicja}

\begin{twierdzenie}[Dołączenie uchwytu Morse'a]
\label{thm:uchwyt-morse-17}
W powyższej sytuacji $M^b$ powstaje,
z dokładnością do zaokrąglenia naroży
i dyfeomorfizmu, z $M^a$ przez
dołączenie uchwytu $H_\lambda$.
W szczególności istnieje odwzorowanie
$S^{\lambda-1}\to M^a$ takie, że
\begin{equation}\label{eq:komorka-morse-17}
 M^b\simeq M^a\cup_{S^{\lambda-1}}D^\lambda,
 \qquad
 H_j(M^b,M^a;\Z)\cong
 \begin{cases}\Z,&j=\lambda,\\0,&j\ne\lambda.
 \end{cases}
\end{equation}
Przy $\lambda=0$ pojawia się składowa;
przy $\lambda=n$ dołącza się
$n$-wymiarowy dysk wzdłuż całej
jego sfery brzegowej.
\end{twierdzenie}
\begin{proof}
\emph{Krok 1: ograniczamy pas.}
Na częściach $[a,c-\varepsilon]$
i $[c+\varepsilon,b]$ nie ma
punktów krytycznych.
Przepływ pola $\nabla f/|\nabla f|^2$
z Twierdzenia~\ref{thm:pas-bez-krytycznych-15}
identyfikuje te pasy z iloczynami.
Wystarczy więc zrozumieć
$f^{-1}([c-\varepsilon,c+\varepsilon])$
dla małego $\varepsilon>0$.

\emph{Krok 2: uchwyt z właściwą częścią przyczepianą.}
W mapie Morse'a mamy $f-c=-|u|^2+|v|^2$.
Wybierzmy $\varepsilon>0$ i $0<r^2<\varepsilon$
tak małe, by poniższy zbiór mieścił się w mapie:
\[
 H=\{(u,v): |v|\leq r,\quad |u|^2\leq\varepsilon+|v|^2\}.
\]
Odwzorowanie
\[
 (a,b)\longmapsto
 \bigl(\sqrt{\varepsilon+r^2|b|^2}\,a,rb\bigr)
 \quad (|a|\leq1,\ |b|\leq1)
\]
jest dyfeomorfizmem $D^\lambda\times D^{n-\lambda}$ na $H$.
Ponadto $H\subset\operatorname{int}M^{c+\varepsilon}$ oraz
\[
 H\cap M^{c-\varepsilon}
 =\{|u|^2=\varepsilon+|v|^2,\ |v|\leq r\}.
\]
Jest to dokładnie obraz $S^{\lambda-1}\times D^{n-\lambda}$;
przy $\lambda=0$ przecięcie jest puste.
Rdzeń $v=0$ ma brzeg w poziomicy $c-\varepsilon$.
Tak zdefiniowana część przyczepiana już leży na właściwym
poziomie; nie trzeba przesuwać jej z płaskiej ściany pudełka.

\emph{Krok 3: dopełnienie jest kołnierzem.}
Wybierzmy metrykę euklidesową na otoczeniu $H$,
przedłużoną podziałem jedności na $M$.
Dla $X=-\nabla f$ mamy tam
\[
 \dot u=2u,\quad\dot v=-2v,\qquad
 (u(t),v(t))=(e^{2t}u(0),e^{-2t}v(0)).
\]
Niech $L=M^{c-\varepsilon}\cup H$.
Blisko naroża jego dopełnienie opisują nierówności
$F=f-(c-\varepsilon)>0$ i $G=|v|^2-r^2>0$.
Na obu ścianach pole jest skierowane do $L$, ponieważ
\[
 XF=-4(|u|^2+|v|^2)<0,\qquad XG=-4|v|^2<0.
\]
Możemy zatem zaokrąglić naroże, zachowując poprzeczność
i ten sam kierunek pola: w płaszczyźnie $(F,G)$
zastępujemy narożnik gładkim łukiem o normalnej
będącej dodatnią kombinacją normalnych obu ścian.
Zaokrąglenie wybieramy na zewnątrz $L$, małe i wewnątrz
$M^{c+\varepsilon}$, tak by zachować $M^{c-\varepsilon}\subset L'$.
Otrzymaną rozmaitość oznaczmy przez $L'$.
Punkt $p$ leży w jej wnętrzu.

Zbiór $K=M^{c+\varepsilon}\setminus\operatorname{int}L'$
jest zwarty i nie zawiera punktów krytycznych.
Stąd $-Xf\geq m>0$ na $K$.
Każda trajektoria z górnego brzegu dociera do $\partial L'$
w skończonym czasie: inaczej $f$ malałaby co najmniej
z prędkością $m$ w zbiorze, na którym jest ograniczona.
W odwrotnym czasie każda trajektoria z $\partial L'$
dociera do górnego brzegu. Nie może powtórnie przeciąć
$\partial L'$, bo $X$ jest tam ściśle skierowane do środka.
Czas trafienia jest gładki na mocy twierdzenia
o funkcji uwikłanej i poprzeczności.
Przeskalowanie czasu na każdej trajektorii daje
$K\cong\partial L'\times[0,1]$.
Wchłonięcie tego kołnierza dowodzi tezy o uchwycie.
Przy $\lambda=n$ nie ma ściany $|v|=r$;
cały brzeg $H=D^n$ przykleja się do dolnej podpoziomicy.
Zamknięte składowe $L'$ nie wymagają żadnego kołnierza.

\emph{Krok 4: komórka i homologia względna.}
Oznaczmy odwzorowanie przyczepienia przez $\alpha(a,b)$.
Homotopia $\alpha(a,(1-t)b)$ zastępuje je odwzorowaniem
zależnym tylko od $a\in S^{\lambda-1}$.
Dołączenia po homotopijnych odwzorowaniach mają ten sam typ
homotopijny pary: homotopię realizuje walec na części
przyczepianej, który wchłaniamy w jej kołnierz.
Po tej zmianie skurczenie czynnika $D^{n-\lambda}$
daje dołączenie samego rdzenia $D^\lambda$.
Nie jest to skurczenie pierwotnego uchwytu
przy jednoczesnym punktowym unieruchomieniu jego całej
części przyczepianej.
Kołnierze pozwalają zastosować wycięcie
(najpierw pogrubiamy część przyczepianą), skąd
\[
 H_j(M^b,M^a;\Z)
 \cong H_j(H_\lambda,S^{\lambda-1}\times D^{n-\lambda};\Z)
 \cong H_j(D^\lambda,S^{\lambda-1};\Z).
\]
Homologia tej pary jest $\Z$ w stopniu $\lambda$
i zerem w pozostałych stopniach, także dla
$(D^0,S^{-1})=(\{\mathrm{punkt}\},\varnothing)$.
\end{proof}

\begin{figure}[htbp]
\centering
\begin{tikzpicture}[>=Latex,font=\small,line cap=round,line join=round]
 \tikzset{old/.style={draw=manifoldblue!80!black,thick,fill=manifoldblue!22},
          new/.style={draw=manifoldgold!85!black,thick,fill=manifoldgold!38},
          seam/.style={manifoldred!80!black,thick,densely dashed},
          pair/.pic={
           \path[old] (-.4,.45)--(-.9,.45)
             arc[start angle=90,end angle=270,radius=.45]
             --(-.4,-.45)--cycle;
           \path[old] (.4,-.45)--(.9,-.45)
             arc[start angle=270,end angle=450,radius=.45]
             --(.4,.45)--cycle;
          },
          ring/.pic={
           \path[old,even odd rule] (0,0) circle (.8) (0,0) circle (.36);
          }}
 \node[font=\bfseries] at (-2.1,3.1) {przed};
 \node[font=\bfseries] at (2.1,3.1) {po};
 \foreach \y in {2.05,0,-2.25}
    \draw[->,gray!70,thick] (-.35,\y)--(.35,\y);
 \node[anchor=east,align=right] at (-3.65,2.05)
    {$0$-uchwyt\\nowa składowa};
 \node[gray!65,font=\Large] at (-2.1,2.05) {$\varnothing$};
 \path[new] (2.1,2.05) circle (.66);
 \draw[gray!22] (-5.1,1.05)--(3.6,1.05);
 \node[anchor=east,align=right] at (-3.65,0)
    {$1$-uchwyt\\połączenie składowych};
 \path (-2.1,0) pic {pair};
 \path (2.1,0) pic {pair};
 \path[new] (1.7,-.45) rectangle (2.5,.45);
 \draw[seam] (1.7,-.45)--(1.7,.45) (2.5,-.45)--(2.5,.45);
 \node[manifoldgold!75!black] at (2.1,-.77) {$D^1\times D^1$};
 \draw[gray!22] (-5.1,-1.2)--(3.6,-1.2);
 \node[anchor=east,align=right] at (-3.65,-2.25)
    {$2$-uchwyt\\wypełnienie otworu};
 \path (-2.1,-2.25) pic {ring};
 \path (2.1,-2.25) pic {ring};
 \path[new] (2.1,-2.25) circle (.36);
 \draw[seam] (2.1,-2.25) circle (.36);
\end{tikzpicture}
\caption{Trzy niezależne przykłady dla dwuwymiarowej
podpoziomicy, pokazane z góry. Niebieski oznacza
część istniejącą, złoty nowy uchwyt, a czerwone
przerywane odcinki lub okrąg --- część przyczepianą.
W drugim wierszu prostokąt przykleja się dokładnie
do dwóch odcinków brzegu. Inne przyklejenie $1$-uchwytu
może utworzyć otwór zamiast połączyć składowe.}
\label{fig:uchwyty-morse-17}
\end{figure}

\begin{przyklad}[Wysokość na sferze]
\label{prz:morse-sfera-17}
Na jednostkowej sferze $S^2\subset\R^3$
weźmy $f(x,y,z)=z$. Krytyczne są
tylko biegun południowy i północny.
Przy południowym
$z=-\sqrt{1-x^2-y^2}
 =-1+\tfrac12(x^2+y^2)+O(|(x,y)|^4)$,
więc indeks wynosi $0$.
Przy północnym znak formy
kwadratowej jest przeciwny,
więc indeks wynosi $2$.
Najpierw pojawia się dysk,
między biegunami podpoziomice
są dyskami, a po przejściu przez
biegun północny dołączamy
$2$-uchwyt. Dostajemy
$c_0=c_2=1$, $c_1=0$.
\end{przyklad}

\begin{przyklad}[Lekko przechylona wysokość torusa]
\label{prz:morse-torus-17}
Standardowy torus parametryzujemy przez
\[
 ((R+r\cos\theta)\cos\phi,\
  (R+r\cos\theta)\sin\phi,\
  r\sin\theta),\qquad R>r>0.
\]
Sama współrzędna $z=r\sin\theta$
ma \emph{okręgi} punktów krytycznych:
nie jest funkcją Morse'a.
Weźmy zamiast niej
\[
 f=z+\varepsilon x
 =r\sin\theta+
    \varepsilon(R+r\cos\theta)\cos\phi,
 \qquad 0<\varepsilon<r/R.
\]
Warunek $\partial_\phi f=0$ daje
$\phi=0$ lub $\phi=\pi$, bo
$R+r\cos\theta>0$.
Warunek $\partial_\theta f=0$ ma
dwa rozwiązania przy każdej z tych
wartości $\phi$, jedno przy górze
i jedno przy dole torusa.
Można je zapisać jawnie. Niech $s=\cos\phi\in\{-1,1\}$,
$t=\operatorname{sgn}(\sin\theta)\in\{-1,1\}$ oraz
$L=\sqrt{1+\varepsilon^2}$ i $h=r/L$.
Równanie $\cos\theta=\varepsilon s\sin\theta$ daje
\[
 \sin\theta=\frac{t}{L},\qquad
 \cos\theta=\frac{\varepsilon st}{L},\qquad
 (x,y,z)=(sR+t\varepsilon h,\,0,\,th).
\]
Wartość funkcji w tym punkcie wynosi
$f=\varepsilon sR+trL$. Te same współrzędne wykorzystano
do umieszczenia znaczników na Rysunku~\ref{fig:torus-morse-17}.
W tych punktach mieszana pochodna
wynosi zero, a na przekątnej hesjanu
stoją
\[
 f_{\theta\theta}
  =-r(\sin\theta+
       \varepsilon\cos\theta\cos\phi),
 \qquad
 f_{\phi\phi}
  =-\varepsilon(R+r\cos\theta)\cos\phi.
\]
Pierwsza jest ujemna u góry
i dodatnia u dołu; druga jest
ujemna dla $\phi=0$
i dodatnia dla $\phi=\pi$.
Zatem indeksy czterech punktów
to kolejno $0,1,1,2$.
Ich wartości, od najmniejszej
do największej, wynoszą
\[
 -r\sqrt{1+\varepsilon^2}-\varepsilon R,\quad
 -r\sqrt{1+\varepsilon^2}+\varepsilon R,\quad
 r\sqrt{1+\varepsilon^2}-\varepsilon R,\quad
 r\sqrt{1+\varepsilon^2}+\varepsilon R.
\]
Nierówność $\varepsilon R<r$
zapewnia tę kolejność.
Podpoziomica jest kolejno:
dyskiem, pierścieniem,
torusem z usuniętym dyskiem
i całym torusem. To są
dołączenia uchwytów indeksów
$0,1,1,2$. Sposób przyklejenia
drugiego $1$-uchwytu rozstrzyga,
że powstaje rodzaj $1$.
\end{przyklad}

\begin{figure}[htbp]
\centering
\begin{tikzpicture}[>=Latex,font=\small,line cap=round]
 \node[font=\bfseries] at (0,1.94)
   {$f(x,y,z)=z+\varepsilon x$, \quad $\varepsilon>0$};
  % Przykladowe epsilon spelnia epsilon R<r; punkty liczymy z df=0.
  \pgfmathsetmacro{\morseeps}{0.16}
  \pgfmathsetmacro{\morseh}{0.69/sqrt(1+\morseeps*\morseeps)}
  \pgfmathsetmacro{\morserplus}{2+\morseeps*\morseh}
  \pgfmathsetmacro{\morserminus}{2-\morseeps*\morseh}
  \begin{axis}[hide axis,view={0}{34},axis equal image,clip=false,scale=2,
    xmin=-3,xmax=3,ymin=-3,ymax=3,zmin=-.8,zmax=.8,scale only axis,
    at={(0,0)},anchor=center,width=9cm,height=5cm,
     colormap={toruslight}{color(0cm)=(manifoldblue!67);
       color(.42cm)=(manifoldblue!33);color(1cm)=(white)}]
   \addplot3[surf,shader=interp,draw=none,z buffer=sort,
     samples=34,samples y=68,domain=0:360,y domain=0:360,
     variable=\t,variable y=\p]
    ({(2+0.69*cos(\t))*cos(\p)},
     {(2+0.69*cos(\t))*sin(\p)},
     {0.69*sin(\t)});
    \addplot3[only marks,mark=*,mark size=2pt,
      mark options={fill=manifoldgreen!75!black,draw=manifoldgreen!75!black}]
      coordinates {(-\morserplus,0,-\morseh)};
    \addplot3[only marks,mark=*,mark size=2pt,
      mark options={fill=manifoldgold!90!black,draw=manifoldgold!90!black}]
      coordinates {(-\morserminus,0,\morseh)};
    \addplot3[only marks,mark=*,mark size=2pt,
      mark options={fill=manifoldgold!90!black,draw=manifoldgold!90!black}]
      coordinates {(\morserminus,0,-\morseh)};
    \addplot3[only marks,mark=*,mark size=2pt,
      mark options={fill=manifoldred!88!black,draw=manifoldred!88!black}]
      coordinates {(\morserplus,0,\morseh)};
    % Punkty, prowadnice i etykiety korzystają z tej samej projekcji.
    \draw[manifoldgreen!70!black,thin]
      (axis cs:-\morserplus,0,-\morseh)--(axis cs:-3.1,0,-1.5)
      node[anchor=east,manifoldgreen!65!black] {minimum $0$};
    \draw[manifoldgold!78!black,thin]
      (axis cs:-\morserminus,0,\morseh)--(axis cs:-3.1,0,1.5)
      node[anchor=east,manifoldgold!75!black] {siodło $1$};
    \draw[manifoldgold!78!black,thin]
      (axis cs:\morserminus,0,-\morseh)--(axis cs:3.1,0,-1.5)
      node[anchor=west,manifoldgold!75!black] {siodło $1$};
    \draw[manifoldred!78!black,thin]
      (axis cs:\morserplus,0,\morseh)--(axis cs:3.1,0,1.5)
      node[anchor=west,manifoldred!78!black] {maksimum $2$};
  \end{axis}
\end{tikzpicture}
\caption{Torus o przekroju kołowym oglądany ukośnie z góry
($R=2$, $r=0{,}69$, $\varepsilon=0{,}16$).
Znaczniki pokazują rzuty wszystkich czterech punktów krytycznych,
także gdy powierzchnia torusa je przesłania. Punkty te mają
indeksy $0,1,1,2$; składnik $\varepsilon x$
rozdziela ich wartości. Podpoziomice są kolejno
dyskiem, pierścieniem, torusem bez dysku i całym
torusem.}
\label{fig:torus-morse-17}
\end{figure}

\section{Komórki, nierówności Morse'a i twierdzenie Reeba}
\label{sec:nierownosci-morse-17}

Przykłady sfery i torusa pokazują, że
punkt krytyczny indeksu $\lambda$ może
dodać tylko \emph{jedną} komórkę wymiaru
$\lambda$. Nie znaczy to, że dodaje
niezależną klasę homologii:
komórka może zabić wcześniejszy cykl.
To właśnie mierzy jej brzeg w kompleksie
komórkowym z
Rozdziału~\ref{ch:homologia-przestrzeni}.

\begin{twierdzenie}[Model CW funkcji Morse'a]
\label{thm:cw-morse-17}
Rozmaitość $M$ ma typ homotopijny
skończonego kompleksu CW $K_f$
z dokładnie $c_k$ komórkami
wymiaru $k$. W szczególności
istnieje kompleks wolnych grup
abelowych
\[
 0\longrightarrow C_n
 \xrightarrow{\partial_n}C_{n-1}
 \longrightarrow\cdots
 \xrightarrow{\partial_1}C_0
 \longrightarrow0,\qquad
 C_k\cong\Z^{c_k},
\]
którego homologiami są $H_k(M;\Z)$.
Generator $C_k$ możemy oznaczyć
punktem krytycznym indeksu $k$.
\end{twierdzenie}
\begin{proof}
Na zwartej $M$ punktów krytycznych
jest skończenie wiele
(Twierdzenie~\ref{thm:gestosc-morse}).
Po małej perturbacji wartości
punktów krytycznych są różne;
wybieramy perturbację stałą
w ich małych otoczeniach,
więc punkty i ich indeksy
pozostają te same, a poza
otoczeniami różniczka $f$
jest oddzielona od zera.
Wybieramy wartości
regularne
$a_0<a_1<\cdots<a_N$
tak, żeby między sąsiednimi
znajdował się dokładnie jeden
punkt krytyczny, a $M^{a_0}$
było puste i $M^{a_N}=M$.
Przy każdym przejściu
Twierdzenie~\ref{thm:uchwyt-morse-17}
zastępuje parę podpoziomic
dołączeniem jednej komórki
$D^{\lambda_i}$.

Wyjaśnijmy indukcję, ponieważ
równoważność homotopijna na
jednym etapie nie jest dosłownie
równością przestrzeni.
Jeśli zbudowany dotąd kompleks
$K_{i-1}$ jest homotopijnie
równoważny $M^{a_{i-1}}$,
przenosimy odwzorowanie
$S^{\lambda_i-1}\to M^{a_{i-1}}$
przez wybraną odwrotność
homotopijną do $K_{i-1}$.
Już na tym etapie stosujemy aproksymację komórkową:
odwzorowanie ze sfery $S^{\lambda_i-1}$ zastępujemy homotopijnym
odwzorowaniem do $(\lambda_i-1)$-szkieletu $K_{i-1}$.
Dzięki temu następne dołączenie rzeczywiście daje
kompleks CW, nawet jeśli starsze komórki mają większy wymiar.
Dla $\lambda_i=0$ po prostu dodajemy punkt.
Dołączenie komórki po odwzorowaniach
homotopijnych daje
homotopijnie równoważne
przestrzenie: homotopię
na sferze rozszerza się
na dołączony walec
$S^{\lambda_i-1}\times[0,1]$,
a następnie wchłania ten
walec w kołnierz dysku.
Tak powstaje $K_i$ z jedną
nową komórką.
Po ostatnim kroku
$K_N\simeq M$.
Złożenie komórek daje
kompleks łańcuchów CW;
ich liczby w każdym stopniu
to dokładnie $c_k$.
\end{proof}

\begin{twierdzenie}[Słabe i mocne nierówności Morse'a]
\label{thm:nierownosci-morse-17}
Niech $b_k=\dim_{\mathbb F}H_k(M;\mathbb F)$
dla dowolnego ciała $\mathbb F$.
Wtedy $c_k\geq b_k$ oraz dla
$0\leq q<n$
\begin{equation}\label{eq:mocne-morse-17}
 \sum_{k=0}^q(-1)^{q-k}c_k
 \ \geq\
 \sum_{k=0}^q(-1)^{q-k}b_k.
\end{equation}
Na końcu otrzymujemy równość
\begin{equation}\label{eq:euler-morse-17}
 \chi(M)=\sum_{k=0}^n(-1)^kc_k.
\end{equation}
\end{twierdzenie}
\begin{proof}
Tensorujemy kompleks z poprzedniego
twierdzenia z $\mathbb F$.
Niech $r_k=\operatorname{rank}_{\mathbb F}
(\partial_k:C_k\otimes\mathbb F
\to C_{k-1}\otimes\mathbb F)$,
przy $r_0=r_{n+1}=0$.
Z twierdzenia o rzędzie i jądrze
\[
 \dim\ker\partial_k=c_k-r_k.
\]
Ponieważ $\partial_k\partial_{k+1}=0$,
obraz $\partial_{k+1}$ jest podprzestrzenią
tego jądra i ma wymiar $r_{k+1}$.
Stąd, \emph{bez pominiętego składnika},
\begin{equation}\label{eq:ranki-morse-17}
 b_k=c_k-r_k-r_{k+1},
 \qquad c_k=b_k+r_k+r_{k+1}.
\end{equation}
Wszystkie rzędy są nieujemne,
więc $c_k\geq b_k$.
Po przemnożeniu~\eqref{eq:ranki-morse-17}
przez $(-1)^{q-k}$ i zsumowaniu
od $0$ do $q$ składniki $r_k$
i $r_{k+1}$ skracają się parami.
Zostaje dokładnie
\[
 \sum_{k=0}^q(-1)^{q-k}(c_k-b_k)
   =r_{q+1}\geq0.
\]
To dowodzi~\eqref{eq:mocne-morse-17}.
Dla $q=n$ prawa strona wynosi
$r_{n+1}=0$. Pomnożenie
przez $(-1)^n$ daje
\eqref{eq:euler-morse-17},
bo $\chi(M)=\sum(-1)^kb_k$.
\end{proof}

\begin{przyklad}[Dolne ograniczenia i torsja]
Na $S^2$ mamy $(b_0,b_1,b_2)=(1,0,1)$;
wysokość z Przykładu~\ref{prz:morse-sfera-17}
realizuje dokładnie dwa niezbędne
punkty. Dla torusa
$(b_0,b_1,b_2)=(1,2,1)$,
a funkcja z
Przykładu~\ref{prz:morse-torus-17}
realizuje $(c_0,c_1,c_2)=(1,2,1)$.
Nierówności możemy stosować
nad różnymi ciałami: dla
$\R P^2$ współczynniki
$\mathbb F_2$ dają
$(b_0,b_1,b_2)=(1,1,1)$,
więc każda funkcja Morse'a
ma co najmniej trzy punkty
krytyczne. Nad $\Q$ środkowa
i najwyższa grupa znikają.
Różnica pokazuje, dlaczego
wybór współczynników może
ujawniać torsję, wspomnianą
już w
Rozdziale~\ref{ch:homologia-przestrzeni}.
\end{przyklad}

\begin{twierdzenie}[Twierdzenie Reeba o dwóch punktach]
\label{thm:reeb-morse-17}
Jeśli spójna zamknięta rozmaitość
$M^n$, $n\geq1$, ma funkcję
Morse'a z dokładnie dwoma
punktami krytycznymi, to
$M$ jest \emph{homeomorficzna}
z $S^n$.
\end{twierdzenie}
\begin{proof}
Minimum i maksimum istnieją,
więc są to oba punkty krytyczne;
ich indeksy wynoszą $0$ i $n$.
Lemat Morse'a daje małą
podpoziomicę wokół minimum
dyfeomorficzną z $D^n$ oraz
małą nadpoziomicę wokół maksimum
również dyfeomorficzną z $D^n$.
Pomiędzy nimi nie ma wartości
krytycznych, więc
Twierdzenie~\ref{thm:pas-bez-krytycznych-15}
przedstawia pas jako
$S^{n-1}\times[0,1]$.
Po wchłonięciu kołnierza
otrzymujemy dwie kopie $D^n$
sklejone wzdłuż brzegu
pewnym dyfeomorfizmem
$h:S^{n-1}\to S^{n-1}$.
Jako \emph{homeomorfizm}
$h$ przedłuża się radialnie
na dysk:
$H(ru)=r\,h(u)$ dla
$u\in S^{n-1}$, $r\in[0,1]$,
oraz $H(0)=0$.
To przedłużenie jest ciągłe
również w $0$, ma ciągłą
odwrotność zbudowaną z $h^{-1}$.
Zmieniając współrzędne
na jednym dysku przez $H$,
otrzymujemy standardowe
sklejenie dwóch półkul,
czyli $S^n$.
\end{proof}

\begin{uwaga}
Ostatni argument jest topologiczny.
Radialne przedłużenie $h$ nie musi być
gładkie w środku dysku, zatem dowód
nie pozwala zastąpić słowa
„homeomorficzna” przez
„dyfeomorficzna”. To rozróżnienie
będzie ważne przy egzotycznych
strukturach gładkich i chirurgii.
\end{uwaga}

\section{Gradient i rozmaitości trajektorii}
\label{sec:gradient-morse-17}

Metryka $g$ przekształca różniczkę
$\dd f$ w gradient $\nabla_g f$
(Definicja~\ref{def:gradient-dywergencja}).
Nie wybieramy jej po to, by zmieniać
punkty krytyczne: te zależą od $f$.
Metryka decyduje natomiast, którędy
biegną trajektorie pola
$X=-\nabla_g f$.
Na zwartej $M$ jego przepływ
$\Phi_t$ istnieje dla wszystkich
$t\in\R$.

\begin{stwierdzenie}[Spadek funkcji i końce trajektorii]
\label{prop:spadek-morse-17}
Dla każdej trajektorii $\gamma(t)=\Phi_t(x)$
\begin{equation}\label{eq:spadek-morse-17}
 \frac{\dd}{\dd t} f(\gamma(t))
 =-|\nabla_g f(\gamma(t))|_g^2\leq0.
\end{equation}
Jeżeli $\gamma$ nie jest stała, to
ma granice przy $t\to-\infty$
i $t\to+\infty$, będące
punktami krytycznymi.
\end{stwierdzenie}
\begin{proof}
Pierwsza równość to reguła łańcuchowa
i definicja gradientu:
$\dd f(-\nabla f)=-g(\nabla f,\nabla f)$.
Równość zachodzi tylko w punkcie
krytycznym; jeśli trajektoria
do niego dotrze w skończonym czasie,
jednoznaczność równania różniczkowego
sprawia, że była stała.
Ponieważ $f(M)$ jest ograniczone,
funkcja $t\mapsto f(\gamma(t))$
ma granice na obu końcach i
\[
 \int_0^\infty
       |\nabla f(\gamma(t))|^2\,\dd t
 =f(x)-\lim_{t\to\infty}f(\gamma(t))
 <\infty.
\]
Na zwartej $M$ pole i pochodna
funkcji $|\nabla f|^2$ są ograniczone,
więc $t\mapsto|\nabla f(\gamma(t))|^2$
jest jednostajnie ciągła.
Nieujemna jednostajnie ciągła
funkcja o skończonej całce na
$[0,\infty)$ dąży do zera:
gdyby przy nieskończenie wielu
chwilach była większa od $\eta>0$,
jednostajna ciągłość dawałaby
rozłączne przedziały stałej
dodatniej długości, na których
przekracza $\eta/2$, sprzeczne
ze skończonością całki.
Każdy punkt skupienia $\gamma(t)$
jest zatem krytyczny.

Zbiór graniczny jest przecięciem
malejących zwartych, spójnych
zbiorów
$\overline{\{\gamma(t):t\geq T\}}$,
więc jest spójny i niepusty.
Punktów krytycznych funkcji
Morse'a na zwartej $M$ jest
skończenie wiele; jedyny
spójny niepusty podzbiór
takiego zbioru to pojedynczy
punkt. Cała trajektoria
zbiega do niego. Dla
$t\to-\infty$ stosujemy
ten sam argument do $-X$.
\end{proof}

\begin{definicja}[Rozmaitości stabilna i niestabilna]
\index{rozmaitosci stabilna i niestabilna@Rozmaitości stabilna i niestabilna}
\label{def:stable-morse-17}
Dla punktu krytycznego $p$ kładziemy
\[
 W^s(p)=\{x:\lim_{t\to+\infty}\Phi_t(x)=p\},
 \qquad
 W^u(p)=\{x:\lim_{t\to-\infty}\Phi_t(x)=p\}.
\]
Litery $s,u$ pochodzą od
angielskich \emph{stable}, \emph{unstable}.
W kierunku przyszłości trajektorie
na $W^s$ zbiegają do $p$;
trajektorie na $W^u$ wychodzą z $p$
przy biegu czasu do przodu.
\end{definicja}

\begin{twierdzenie}[Wymiary rozmaitości stabilnych]
\label{thm:stable-morse-17}
Dla dowolnej gładkiej metryki $g$ rozmaitość
$W^u(p)$ jest różnowartościowo
immersowaną rozmaitością
dyfeomorficzną z $\R^{\lambda(p)}$,
a $W^s(p)$ z $\R^{n-\lambda(p)}$.
Ich przestrzenie styczne w $p$ są odpowiednio
dodatnią i ujemną przestrzenią spektralną $DX_p$.
Metrykę można ponadto wybrać euklidesową
w małych rozłącznych mapach Morse'a. W tych mapach
\[
 W^u_{\mathrm{loc}}(p)=\{v=0\},
 \qquad
 W^s_{\mathrm{loc}}(p)=\{u=0\}.
\]
\end{twierdzenie}
\begin{proof}
Ponieważ punktów krytycznych
jest skończenie wiele, bierzemy
rozłączne mapy z
Lematu~\ref{thm:lemat-morsea-17}
i gładkie funkcje odcinające
równe $1$ blisko ich środków.
Wypukłe połączenie
metryk euklidesowych w tych
mapach z dowolną metryką
globalną daje metrykę
euklidesową na mniejszych mapach.
Tam równania przepływu
i ich rozwiązania podaliśmy
w dowodzie
Twierdzenia~\ref{thm:uchwyt-morse-17}.
Warunek zbiegania do $(0,0)$
przy $t\to+\infty$ zmusza $u=0$,
bo $u(t)=e^{2t}u(0)$;
analogicznie zbieganie
przy $t\to-\infty$ zmusza $v=0$.
To dowodzi lokalnego opisu
i wymiarów.

Każdy $x\in W^s(p)$
po dostatecznie długim czasie
wpada do małego lokalnego
dysku $\{u=0\}$, a
$x=\Phi_{-T}(y)$
dla pewnego punktu $y$
tego dysku. Stąd
\[
 W^s(p)=\bigcup_{T\geq0}
      \Phi_{-T}(W^s_{\mathrm{loc}}(p)).
\]
Te obrazy są gładkimi kopiami
otwartego dysku i są
zagnieżdżone, jeśli lokalny
dysk wybierzemy niezmienniczy
przy przepływie naprzód.
Odwzorowania przejścia są
ograniczeniami dyfeomorfizmów
$\Phi_t$, więc określają
gładką strukturę na sumie.
Ustalmy w lokalnym dysku
kulę promienia $\delta$.
Każda trajektoria w
$W^s(p)\setminus\{p\}$
przecina jej sferę raz, gdy
wbiega do tej kuli: od tej
chwili lokalny wzór
$v(t)=e^{-2t}v(0)$
ściśle zmniejsza $|v|$.
Punkt takiej trajektorii
opisują kierunek na
$S^{n-\lambda(p)-1}$
i czas od przecięcia.
Jawna parametryzacja, dla $v\ne0$, to
\[
 v\longmapsto
 \Phi_{-\frac12\log(|v|/\delta)}
          (0,\delta v/|v|).
\]
Blisko zera jest to dokładnie $(0,v)$,
więc przedłuża się gładko przez $v=0$.
Jedyność przecięcia sfery i gładkość czasu trafienia
dają odwrotność, a więc standardową
strukturę $\R^{n-\lambda(p)}$.
Włączenie w $M$ jest
różnowartościową immersją.
Dla wymiaru zerowego
jest to po prostu $\{p\}$.
Dla $W^u$ odwracamy czas;
wynik ma wymiar $\lambda(p)$.
Nie twierdzimy, że globalna
immersja jest właściwa.

Uzasadnijmy teraz twierdzenie dla ogólnej metryki.
W punkcie krytycznym
$A=DX_p=-g_p^{-1}\operatorname{Hess}_p f$ jest
samosprzężony względem $g_p$ i nie ma zerowych
wartości własnych. Ma $\lambda(p)$ dodatnich
i $n-\lambda(p)$ ujemnych wartości własnych.
W liniowych współrzędnych spektralnych równanie ma postać
$\dot x=Ax+R(x)$, gdzie $R(0)=DR(0)=0$.
Oznaczmy części stabilną i niestabilną przez $s,u$.
W odpowiedniej normie, dla pewnego $\alpha>0$,
\[
 \|e^{A_st}\|\leq e^{-\alpha t},\qquad
 \|e^{-A_ut}\|\leq e^{-\alpha t}\quad(t\geq0).
\]
Po zmniejszeniu kuli i odcięciu poza nią
stała Lipschitza $R$ może być dowolnie mała.
Dla zadanego $x_s(0)=a$ rozwiązanie zbiegające
do zera musi spełniać
\begin{align*}
 x_s(t)&=e^{A_st}a+
       \int_0^t e^{A_s(t-\tau)}R_s(x(\tau))\,\dd\tau,\\
 x_u(t)&=-\int_t^\infty
       e^{A_u(t-\tau)}R_u(x(\tau))\,\dd\tau.
\end{align*}
Drugi wzór wynika z wariacji stałych i znikania
rosnącego rozwiązania jednorodnego w nieskończoności.
W przestrzeni funkcji z normą
$\|x\|_\beta=\sup_{t\geq0}e^{\beta t}|x(t)|$,
gdzie $0<\beta<\alpha$, prawa strona jest kontrakcją:
jej stała jest nie większa niż
\[
 \operatorname{Lip}(R)
 \left((\alpha-\beta)^{-1}+(\alpha+\beta)^{-1}\right)<1.
\]
Twierdzenie o punkcie stałym daje jedyne rozwiązanie.
Gładka zależność od $a$ wynika z różniczkowania
tego równania: operator $I-D\mathcal T$ ma odwrotność
$\sum_{j\geq0}(D\mathcal T)^j$; argument powtarzamy
dla kolejnych pochodnych. Otrzymujemy gładki wykres
$x_u(0)=h(a)$, z $h(0)=Dh(0)=0$.
Każda trajektoria pozostająca w małej kuli na przyszłość
i zbiegająca do zera leży na tym wykresie:
spełnia te same równania całkowe, które są kontrakcją
także w normie supremum ($\beta=0$).
Jest to lokalna rozmaitość stabilna.
Odwrócenie czasu daje niestabilną.

Na wykresie stabilnym, w parametrze $a$, mamy
$\dot a=A_sa+O(|a|^2)$, zatem
$\frac{\dd}{\dd t}|a|^2<0$ poza zerem w małej kuli.
Przepływ jest poprzeczny do jej sfery brzegowej.
Do tej kuli globalna rozmaitość stabilna dokleja
więc zewnętrzny nieskończony kołnierz sfery,
parametryzowany ujemnym czasem. Kula z takim
kołnierzem jest dyfeomorficzna z $\R^{n-\lambda(p)}$:
na kołnierzu zamieniamy czas na promień,
zgodnie z lokalnym kołnierzem brzegu kuli.
Ten sam argument z odwróconym czasem kończy dowód.
\end{proof}

\begin{figure}[htbp]
\centering
\begin{tikzpicture}[>=Latex,scale=1]
 \shade[inner color=white,outer color=manifoldblue!43]
   (0,0) circle (1.72);
 \draw[manifoldblue!70!black,thick] (0,0) circle (1.72);
 % W tym rzucie prostopadłym równik jest odcinkiem.
 \draw[manifoldblue!45,dashed] (-1.72,0)--(1.72,0);
 % Rzuty południków: (R cos(phi) sin(theta), R cos(theta)).
 \foreach \c in {-.85,-.45,0,.45,.85} {
   \draw[manifoldgold!90!black,thick,domain=10:170,samples=65]
      plot ({1.72*\c*sin(\x)},{1.72*cos(\x)});
   \draw[manifoldgold!90!black,thick,->,domain=76:105,samples=20]
      plot ({1.72*\c*sin(\x)},{1.72*cos(\x)});
 }
 \fill[manifoldred] (0,1.72) circle (2.4pt)
    node[above] {$p_+,\ \lambda=2$};
 \fill[manifoldblue!85!black] (0,-1.72) circle (2.4pt)
    node[below] {$p_-,\ \lambda=0$};
 \node[manifoldgold!75!black] at (2.45,.2)
       {$-\nabla f$};
\end{tikzpicture}
\caption{Spadek funkcji wysokości na sferze.
Dla metryki indukowanej trajektorie są południkami
i biegną od maksimum do minimum; rysunek pokazuje ich rzuty.
Południowy punkt jest stabilny w dwóch
wymiarach, a północny niestabilny w dwóch
wymiarach; pozostałe zbiory stabilne
lub niestabilne tych punktów to same punkty.}
\label{fig:przeplyw-morse-sfera-17}
\end{figure}

\section{Warunek Morse'a--Smale'a i przestrzenie trajektorii}
\label{sec:morse-smale-17}

Dwie rozmaitości $W^u(p)$ i $W^s(q)$
mogą spotkać się stycznie. Wtedy
niewielka zmiana metryki zmienia
obraz trajektorii w trudny do
kontrolowania sposób. Warunek
transwersalności usuwa ten problem.

\begin{definicja}[Para Morse'a--Smale'a]
\index{para morsea--smalea@Para Morse'a--Smale'a}
\label{def:ms-17}
Para $(f,g)$ jest \emph{parą Morse'a--Smale'a},
gdy $f$ jest funkcją Morse'a
i dla każdych punktów krytycznych
$p,q$ rozmaitości $W^u(p)$ i $W^s(q)$
przecinają się transwersalnie.
Jeśli $p\ne q$, oznaczamy przez
\[
 \mathcal M(p,q)
 :=(W^u(p)\cap W^s(q))/\R
\]
zbiór trajektorii z $p$ do $q$
z zapomnianym początkiem czasu;
$\R$ działa przez przesunięcie
parametru $t$.
\end{definicja}

\begin{stwierdzenie}[Wymiar trajektorii]
\label{prop:wymiar-ms-17}
Jeśli para jest Morse'a--Smale'a
i $\mathcal M(p,q)$ jest niepusta,
to $\lambda(p)>\lambda(q)$ oraz
\[
 \dim\mathcal M(p,q)
    =\lambda(p)-\lambda(q)-1.
\]
Przy różnicy indeksów równej $1$
trajektorie są izolowane.
\end{stwierdzenie}
\begin{proof}
Transwersalność i
Twierdzenie~\ref{thm:stable-morse-17}
dają
\[
 \dim(W^u(p)\cap W^s(q))
 =\lambda(p)+(n-\lambda(q))-n
 =\lambda(p)-\lambda(q).
\]
W punkcie niekrytycznym przecięcia
wektor $-\nabla f$ jest niezerowy
i styczny do obu rozmaitości,
bo obie są niezmiennicze
przez przepływ. Dlatego
przecięcie ma wymiar co
najmniej $1$; stąd
$\lambda(p)>\lambda(q)$.
Koniecznie $f(p)>f(q)$. Wybierzmy regularne
$s$ między tymi wartościami i oznaczmy
$\Sigma=W^u(p)\cap W^s(q)\cap f^{-1}(s)$.
Każda trajektoria przecina $\Sigma$ dokładnie raz.
Czas tego przecięcia jest gładki, ponieważ
$Xf\ne0$ na poziomicy (twierdzenie o funkcji uwikłanej).
Przepływ daje zatem globalny dyfeomorfizm
\[
 \Sigma\times\R\longrightarrow W^u(p)\cap W^s(q),
 \qquad (z,t)\longmapsto\Phi_t(z).
\]
Przecięcia rozumiemy w strukturach rozmaitości
immersowanych, czyli jako iloczyny włókniste ich
odwzorowań do $M$. Ich różnowartościowość pozwala
utożsamiać punkty z punktami przecięcia w $M$.
W konsekwencji iloraz jest rozmaitością Hausdorffa
dyfeomorficzną z $\Sigma$, o wymiarze mniejszym o jeden.
Samo lokalne odjęcie wymiaru działania nie
wystarczałoby do uzasadnienia tych własności ilorazu.
\end{proof}

\begin{twierdzenie}[Istnienie metryk Morse'a--Smale'a]
\label{thm:gen-ms-17}
Dla ustalonej funkcji Morse'a
na zwartej $M$ można wybrać
metrykę $g$, euklidesową w małych
mapach Morse'a, tak aby
$(f,g)$ była parą
Morse'a--Smale'a. Takie metryki
są gęste wśród metryk
z ustalonym zachowaniem
blisko punktów krytycznych.
\end{twierdzenie}
\begin{proof}
Pokażemy mechanizm
transwersalności, a następnie
zastosujemy
Twierdzenie~\ref{thm:transwersalnosc-param}.
Ustalmy różne punkty krytyczne
$p,q$ z $f(p)>f(q)$ i regularny
poziom $L=f^{-1}(s)$ między nimi.
Ślady $W^u(p)\cap L$ i $W^s(q)\cap L$
są gładkimi rozmaitościami immersowanymi w $L$.
Wybieramy $s$ poza ustalonymi małymi otoczeniami
punktów krytycznych; poniższe perturbacje
nie będą więc zmieniały metryki w tych otoczeniach.
Wystarczy uzyskać ich
transwersalność na $L$.

Na obszarze, gdzie $\dd f\ne0$,
zmiana metryki może dowolnie
zmienić gradient w pierwszym
rzędzie w kierunkach poprzecznych
do trajektorii. Istotnie, dla
$X=-g^{-1}\dd f$ i symetrycznej
perturbacji $h$ metryki mamy
\[
 \left.\frac{\dd}{\dd\varepsilon}
  (- (g+\varepsilon h)^{-1}\dd f)
 \right|_{\varepsilon=0}
  =-g^{-1}hX.
\]
Gdy $X\ne0$, odwzorowanie
$h\mapsto h(X,\cdot)$ z form
symetrycznych na kowektory
jest surjektywne: w bazie
z pierwszym wektorem
równoległym do $X$ możemy
zadać dowolny pierwszy
wiersz i tę samą pierwszą
kolumnę macierzy $h$.
Wybieramy kulkę na tej
trajektorii \emph{tuż ponad} $L$
i perturbację podpartą tylko
w tej kulce. Poniżej $L$
metryka pozostaje taka sama,
więc $W^s(q)\cap L$ nie zmienia
się. Zmiana wektora prędkości
w kierunku stycznym do $L$
zmienia punkt trafienia w $L$
w tym samym kierunku
w pierwszym rzędzie: wynika
to z całkowego równania
dla zlinearyzowanego przepływu;
przy zmniejszaniu kulki
przeniesienie liniowe na
krótkim odcinku zbliża się
do identyczności.
Wpływ perturbacji na
trajektorię zależy gładko
od jej skończenie wielu
parametrów na mocy
gładkiej zależności
rozwiązań ODE od parametrów.

Doprecyzujmy argument gęstości, ponieważ same zbiory
$W^u,W^s$ zależą od metryki. Ustalmy małą sferę
wyjściową $S_p\subset W^u_{\mathrm{loc}}(p)$
i wejściową $S_q\subset W^s_{\mathrm{loc}}(q)$.
Są one niezależne od zmienianej metryki.
Niech $E_g(x)$ oznacza pierwsze trafienie trajektorii
z $x\in S_p$ w $L$, a $D_g(y)$ pierwsze trafienie
trajektorii biegnącej wstecz z $y\in S_q$ w $L$.
Odwzorowania są określone na otwartych zbiorach tych punktów,
które osiągają $L$ w skończonym czasie.
Zależą gładko od punktu i od skończenie wielu
parametrów metryki. Szukany warunek to
transwersalność $(E_g,D_g)$ do przekątnej w $L\times L$.

Dla każdego $N\in\mathbb N$ żądamy transwersalności
we wszystkich przecięciach, których oba czasy
trafienia są nie większe niż $N$.
Jest to warunek otwarty: ewentualny ciąg przecięć
naruszających go przy $g_j\to g$ ma podciąg
zbieżnych punktów na zwartych $S_p,S_q$
i zbieżnych czasów w $[0,N]$; gładka zależność
przepływu przeczyłaby transwersalności dla $g$.
Jest też gęsty. Przy ustalonej metryce zbiór
takich par startowych jest zwarty, więc wystarcza
skończenie wiele małych obszarów perturbacji
opisanych wyżej. Pochodne parametryczne pierwszej
ewaluacji rozpinają $TL$, podczas gdy druga
ewaluacja nie zmienia się. Parametryczna
transwersalność zastosowana w otoczeniu tego
zwartego zbioru daje dowolnie małe perturbacje
spełniające warunek. Nowe przecięcia o czasach
$\leq N$ nie mogą pojawić się poza tym otoczeniem,
co wynika z tego samego argumentu zwartości.

Przestrzeń gładkich metryk o ustalonych wartościach
w otoczeniach krytycznych jest otwartą częścią
zamkniętej afinicznej podprzestrzeni przestrzeni
Fr\'echeta gładkich tensorów symetrycznych.
Ma własność Baire'a. Przecięcie powyższych
otwartych gęstych warunków po wszystkich $N$
jest więc gęste i obejmuje każde przecięcie
o skończonych czasach trafienia.
Małe perturbacje zachowują dodatnią określoność.
Punktów krytycznych jest
skończenie wiele, więc
przeprowadzamy to dla
wszystkich par. Jeśli
$f(p)\leq f(q)$ i $p\ne q$,
trajektoria spadku nie może
biec z $p$ do $q$;
przy $p=q$ lokalne osie
$\{u=0\}$ i $\{v=0\}$
przecinają się transwersalnie
w $p$.
\end{proof}

\begin{figure}[htbp]
\centering
\begin{tikzpicture}[>=Latex,scale=.95]
 \draw[manifoldblue!30,densely dashed]
   (-2.25,.55) .. controls (-.9,.49) and (.9,.61) .. (2.25,.55);
 \draw[manifoldred!80!black,very thick,->]
   (-1.4,1.50) .. controls (-1.23,.55) and (-.48,.1)
                .. (-.1,-.25);
 \draw[manifoldred!80!black,very thick,->]
   (-.1,-.25) .. controls (.13,-.66) and (.58,-1.14)
             .. (.9,-1.41);
 \fill[manifoldred] (-1.4,1.5) circle (2.3pt)
    node[above] {$p$};
 \fill[manifoldgold!85!black] (-.1,-.25) circle (2.3pt)
    node[left] {$r$};
 \fill[manifoldblue!85!black] (.9,-1.41) circle (2.3pt)
    node[below] {$q$};
 \node[manifoldblue!70!black] at (2.65,.55) {$f=s$};
\end{tikzpicture}
\caption{Schemat trajektorii łamanej $p\to r\to q$.
Pojedyncza trajektoria od $p$ do $q$
może się do niej zbliżać, spędzając
coraz więcej czasu blisko $r$.
Taka granica jest punktem brzegowym
jednowymiarowej przestrzeni trajektorii,
gdy różnica indeksów $p,q$ wynosi $2$.}
\label{fig:lamane-trajektorie-morse-17}
\end{figure}

\section{Kompleks Morse'a: homologia z trajektorii}
\label{sec:kompleks-morse-17}

Teraz ustalmy metrykę Morse'a--Smale'a.
Nie mylmy symboli: w \emph{homologii}
stosujemy operator brzegu $\partial$,
zaś $d$ pozostaje dla form
i operatorów kohomologii
z poprzednich rozdziałów.

\begin{lemat}[Zwartość z dopuszczeniem przełamań]
\label{lem:zwartosc-lamana-17}
Z każdego ciągu trajektorii od $p$ do $q$ można wybrać
podciąg zbiegający do skończonego łańcucha trajektorii
$p=r_0\to r_1\to\cdots\to r_\ell=q$.
Zbieżność oznacza gładką zbieżność na zwartych
przedziałach czasu po osobnym przesunięciu czasu
dla każdego niestałego odcinka.
\end{lemat}
\begin{proof}
Wybierzmy regularne poziomy rozdzielające wszystkie
różne wartości krytyczne między $f(q)$ i $f(p)$.
Każda trajektoria przecina każdy z tych poziomów raz.
Zwartość ich skończonego iloczynu daje podciąg,
na którym wszystkie punkty przecięć zbiegają.
Przepływ z każdego punktu granicznego wyznacza
graniczną trajektorię; gładka zależność od danych
początkowych daje wymaganą zbieżność lokalną.
Na zwartym paśmie bez zer $|\nabla f|\geq m>0$,
więc czas jego przejścia nie przekracza
różnicy poziomów podzielonej przez $m^2$.
Tam odcinki graniczne muszą się łączyć.

Nieograniczone czasy mogą zatem występować tylko
w otoczeniach poziomów krytycznych.
Końce granicznych odcinków są punktami krytycznymi
na mocy Stwierdzenia~\ref{prop:spadek-morse-17}.
Trzeba jeszcze wykluczyć, że odcinek przychodzi do
jednego punktu, a następny wychodzi z innego punktu
\emph{tej samej} wartości krytycznej.
Weźmy rozłączne małe kule wokół tych punktów.
Na zwartym obszarze między kulami $|\nabla f|\geq m>0$,
a przejście między różnymi kulami wymaga długości
co najmniej $d>0$. Spadek funkcji na tym przejściu wynosi
\[
 \int |\nabla f|^2\,\dd t
 =\int |\nabla f|\,\dd s\geq md,
\]
gdzie $s$ jest długością łuku. Nie może on zmieścić się
w pasie o szerokości dążącej do zera.
Tak więc oba odcinki mają wspólny punkt krytyczny.
To samo rozumowanie na małych sferach wyjścia
z $p$ i wejścia do $q$ zachowuje końce całego łańcucha.
Każdy niestały odcinek ściśle obniża $f$,
a więc łańcuch ma skończoną długość.

Zapis przez przecięcia z wybranymi poziomami
wyznacza jednoznacznie trajektorię łamaną.
Nadaje zbiorowi takich łańcuchów topologię Hausdorffa
pochodzącą z metryki skończonego iloczynu poziomic.
Powyższy argument stosuje się też do ciągu łańcuchów:
wybieramy podciąg o ustalonej liście punktów
krytycznych i stosujemy go do każdego odcinka.
Domknięcie przestrzeni trajektorii w tym zbiorze,
oznaczane $\overline{\mathcal M}(p,q)$, jest zatem zwarte.
\end{proof}

\begin{lemat}[Zwarta przestrzeń trajektorii indeksu jeden]
\label{lem:fin-morse-17}
Jeśli $\lambda(p)-\lambda(q)=1$,
to $\mathcal M(p,q)$ jest
skończonym zbiorem.
\end{lemat}
\begin{proof}
Przestrzeń jest zerowymiarowa przez
Stwierdzenie~\ref{prop:wymiar-ms-17}.
Każdy odcinek trajektorii łamanej ściśle obniża indeks.
Przy różnicy indeksów $1$ nie ma miejsca na punkt pośredni.
Lemat~\ref{lem:zwartosc-lamana-17} daje więc zwartość
samej $\mathcal M(p,q)$. Zwarta rozmaitość
zerowymiarowa jest skończonym zbiorem.
\end{proof}

Wybierzmy orientację każdej
rozmaitości $W^u(p)$.
Można to zrobić osobno,
bo jest dyfeomorficzna
z przestrzenią euklidesową;
nie wymaga to orientowalności
całego $M$. Dla trajektorii
$\gamma:p\to q$ o różnicy
indeksów $1$ orientacje
obu rozmaitości niestabilnych
i kierunek przepływu określają
znak $\epsilon(\gamma)\in\{\pm1\}$.
Konkretnie w regularnej
poziomicy $L$ jest to znak
poprzecznego przecięcia
$W^u(p)\cap L$ z
$W^s(q)\cap L$;
normalną orientację drugiego czynnika
przenosimy z $T_qW^u(q)$ wzdłuż $W^s(q)$.
Jest to jednoznaczne, ponieważ $W^s(q)$ jest ściągalna.
Orientację $W^u(p)\cap L$ ustalamy regułą
\[
 \operatorname{or}(TW^u(p))
 =\operatorname{or}(X)\otimes
   \operatorname{or}(T(W^u(p)\cap L)).
\]
Dla różnicy indeksów $1$ znak jest zatem znakiem
izomorfizmu z przestrzeni stycznej tego śladu
do zorientowanej przestrzeni normalnej $W^s(q)$.
Przy większej różnicy indeksów orientujemy jądro
tego odwzorowania, czyli $T\mathcal M(p,q)$,
stawiając je przed przestrzenią normalną.
Zmiana orientacji jednego
$W^u(p)$ zmienia znaki
odpowiedniego generatora,
lecz nie grupy homologii.

\begin{definicja}[Kompleks trajektorii Morse'a]
\index{kompleks trajektorii morsea@Kompleks trajektorii Morse'a}
\label{def:kompleks-morse-17}
Grupa $CM_k(f,g;\Z)$ jest wolną
grupą abelową na punktach
$\operatorname{Crit}_k(f)$.
Dla generatora $p$ kładziemy
\begin{equation}\label{eq:brzeg-morse-17}
 \partial^{\mathrm{Mor}}_k p
 :=\sum_{q\in\operatorname{Crit}_{k-1}(f)}
  \left(\sum_{\gamma\in\mathcal M(p,q)}
                    \epsilon(\gamma)\right)q.
\end{equation}
Suma jest skończona przez
Lemat~\ref{lem:fin-morse-17}.
Nad $\mathbb F_2$ zapominamy znaki:
współczynnikiem przy $q$
jest liczba trajektorii modulo $2$.
\end{definicja}

\begin{twierdzenie}[Dlaczego $\partial^2=0$]
\label{thm:brzeg-morse-17}
Operator~\eqref{eq:brzeg-morse-17}
spełnia
$\partial^{\mathrm{Mor}}_{k-1}
 \partial^{\mathrm{Mor}}_k=0$.
\end{twierdzenie}
\begin{proof}
Ustalmy $p$ o indeksie $k$
i $q$ o indeksie $k-2$.
Ze Stwierdzenia~\ref{prop:wymiar-ms-17}
przestrzeń $\mathcal M(p,q)$
jest jednowymiarowa.
Pokażemy, jakie granice
mogą mieć ciągi jej elementów.
Argument ze zwartością
poziomic z
Lematu~\ref{lem:fin-morse-17}
daje podciąg zbieżny
albo trajektorię $p\to q$,
albo trajektorię łamaną.
Każde przełamanie zużywa
co najmniej jednostkę
różnicy indeksów; przy
różnicy $2$ możliwe
jest dokładnie jedno
$p\to r\to q$ z
$\lambda(r)=k-1$.

Wyjaśnijmy, dlaczego każdy taki łańcuch ma
dokładnie jeden kołnierz, a więc jest rzeczywistym końcem.
Oznaczmy $l=\lambda(r)$, $b=n-l$.
Najpierw załóżmy, że metryka w mapie Morse'a jest euklidesowa.
Przekrój wejściowy ma postać $|v|=\delta$, z parametrami
$(u,\theta)$, gdzie $\theta=v/\delta\in S^{b-1}$.
Przekrój wyjściowy $|u|=\delta$ ma parametry
$(\omega,v)$, gdzie $\omega=u/\delta\in S^{l-1}$.
Transwersalność przy ustalonych trajektoriach $p\to r$
i $r\to q$ pozwala zapisać ślady odpowiednich gałęzi jako
\[
 W^u(p)\cap\Sigma_{\rm in}:\quad\theta=A(u),
 \qquad
 W^s(q)\cap\Sigma_{\rm out}:\quad\omega=B(v).
\]
Istotnie, pierwszy ślad ma wymiar $l$ i jest
poprzeczny do $u=0$, a drugi ma wymiar $b$
i jest poprzeczny do $v=0$.
Przy czasie przejścia $T$ i $\rho=e^{-2T}$
lokalne równania przepływu dają
\[
 u_{\rm in}=\delta\rho\omega,\quad
 v_{\rm in}=\delta\theta,\quad
 u_{\rm out}=\delta\omega,\quad
 v_{\rm out}=\delta\rho\theta.
\]
Warunek dopasowania to dokładnie
\begin{equation}\label{eq:sklejanie-morse-17}
 \theta=A(\delta\rho\omega),\qquad
 \omega=B(\delta\rho\theta).
\end{equation}
Przy $\rho=0$ rozwiązaniem jest $(A(0),B(0))$,
a różniczka równań względem współrzędnych kątowych
$(\theta,\omega)$ jest identycznością.
Twierdzenie o funkcji uwikłanej daje zatem
\emph{jedyne} pobliskie rozwiązanie dla każdego małego
$\rho\geq0$. Dla $\rho>0$ jest to trajektoria niełamana,
a $\rho=0$ oznacza zadany łańcuch.
Dla sfer $S^0$ współrzędne kątowe po prostu nie występują.

Ogólna metryka wymaga dodatkowego argumentu;
nie wolno zakładać gładkiej linearyzacji jej gradientu.
Wyprostujmy lokalne rozmaitości stabilną i niestabilną
z Twierdzenia~\ref{thm:stable-morse-17} do osi $u=0,v=0$.
Równania mają teraz postać
$\dot u=A_uu+R_u(u,v)$, $\dot v=A_sv+R_s(u,v)$,
gdzie $R_u(0,v)=R_s(u,0)=0$ i $DR(0)=0$.
W małym pudełku $|u|,|v|\leq\delta$ norma $u$
ściśle rośnie, a norma $v$ ściśle maleje poza odpowiednią osią.
Dla warunków mieszanych $v(0)=a$, $u(T)=b_0$
wariacja stałych daje
\begin{align*}
 u(t)&=e^{A_u(t-T)}b_0-
       \int_t^T e^{A_u(t-\tau)}R_u(u(\tau),v(\tau))\,\dd\tau,\\
 v(t)&=e^{A_st}a+
       \int_0^t e^{A_s(t-\tau)}R_s(u(\tau),v(\tau))\,\dd\tau.
\end{align*}
Tak jak w dowodzie rozmaitości stabilnej,
małość $DR$ czyni ten operator kontrakcją,
jednostajnie względem $T$.
Stosujemy normę ważoną
\[
 \sup_{0\leq t\leq T}
 \bigl(e^{\eta(T-t)}|u(t)|+e^{\eta t}|v(t)|\bigr),
 \qquad 0<\eta<\alpha.
\]
Znikanie $R_u$ na osi stabilnej i $R_s$ na osi
niestabilnej kontroluje wyrazy mieszane: na kuli
tej normy pochodne mieszane są ograniczone przez
stałą razy $\delta e^{-\eta(T-t)}$ lub
$\delta e^{-\eta t}$. Całkowanie jąder wykładniczych
daje stałą kontrakcji mniejszą od $1$ dla małego $\delta$.
Różniczkowanie równania punktu stałego daje
te same oszacowania dla pochodnych po $a,b_0$.
W szczególności odwzorowania
\[
 U_T(a,b_0)=u(0),\qquad V_T(a,b_0)=v(T)
\]
dążą do zera w normie $C^1$ z szybkością
$O(e^{-\eta T})$ na małych ustalonych zbiorach parametrów.
Zależą gładko od $T<\infty$; można to zobaczyć,
przenosząc równania całkowe na ustalony przedział $[0,1]$.

Równania dopasowania zastępujemy teraz przez
\[
 \theta=A(U_T(\delta\theta,\delta\omega)),\qquad
 \omega=B(V_T(\delta\theta,\delta\omega)).
\]
Dla dużego $T$ ich prawa strona jest kontrakcją
w małym iloczynie map wokół $(A(0),B(0))$.
Istnieje więc dokładnie jedno rozwiązanie,
gładkie dla skończonego $T$ i zbieżne do rozwiązania
łamanego dla $T\to\infty$.
Parametr $1/T\geq0$ daje topologiczny kołnierz końca,
gładki w jego wnętrzu. Możemy go przyjąć za mapę
gładkiej struktury rozmaitości z brzegiem.
Nie potrzebujemy tu gładkości odwzorowania ewaluacji
w zmiennej $1/T$ na samym brzegu.
Jedyność rozwiązania i monotoniczność norm zapewniają,
że wszystkie dostatecznie bliskie trajektorie
są opisane przez ten kołnierz.
Wspomniany argument
zwartości wyklucza inne
końce. Dodając je,
otrzymujemy zwartą
rozmaitość jednowymiarową
z brzegiem
\[
 \partial\overline{\mathcal M}(p,q)
  =\coprod_{\lambda(r)=k-1}
      \mathcal M(p,r)\times\mathcal M(r,q).
\]
Orientacje ustalone przed
Definicją~\ref{def:kompleks-morse-17} orientują wnętrze
tej rozmaitości. Porównanie z orientacją końca można
przeprowadzić na powyższych równaniach dopasowania.
Parametryzacje obu wykresów wnoszą znaki
$\epsilon(\gamma_{pr})$ i $\epsilon(\gamma_{rq})$,
a macierz dopasowania dąży do identyczności,
więc jej wyznacznik jest dodatni.
Pozostaje przestawienie bloków: kierunku przepływu,
parametru kołnierza i przestrzeni normalnych.
Jego znak zależy tylko od wymiarów $k,n$, nie od
trajektorii ani od $r$ (wszystkie te $r$ mają indeks $k-1$).
Dlatego znak punktu brzegowego jest
\[
 \sigma_{k,n}\,
 \epsilon(\gamma_{pr})\epsilon(\gamma_{rq}),
 \qquad \sigma_{k,n}\in\{1,-1\},
\]
z jednym wspólnym czynnikiem konwencji orientacyjnej.
W przypadku ogólnej metryki macierz dopasowania
również dąży do identyczności, więc porównanie jest to samo.
Każdy zwarty zorientowany odcinek ma dwa końce
o przeciwnych znakach, a okrąg nie ma brzegu.
Suma znaków brzegu jest zerem; po pomnożeniu przez
$\sigma_{k,n}$ dostajemy
\[
 \sum_{\lambda(r)=k-1}
 \sum_{\gamma_{pr},\gamma_{rq}}
 \epsilon(\gamma_{pr})\epsilon(\gamma_{rq})=0.
\]
Jest to dokładnie współczynnik przy $q$ w
$\partial^{\mathrm{Mor}}\partial^{\mathrm{Mor}}p$.
Po redukcji modulo $2$
to samo rozumowanie
mówi, że liczba
końców jest parzysta.
\end{proof}

\begin{twierdzenie}[Homologia Morse'a jest homologią $M$]
\label{thm:homologia-morse-17}
Dla zamkniętej rozmaitości
i pary Morse'a--Smale'a
\[
 H_k(CM_\bullet(f,g;\Z),
          \partial^{\mathrm{Mor}})
       \cong H_k(M;\Z).
\]
Wynik nie zależy od
wybranego $f$, metryki
ani orientacji generatorów.
\end{twierdzenie}
\begin{proof}
\emph{Krok 1: filtracja według indeksu.}
Sama zgodność liczby generatorów z liczbą komórek
nie porównuje różniczek. Potrzebna jest ponadto
filtracja uporządkowana według indeksów.
Użyjemy geometrycznego twierdzenia o przestawianiu
wartości krytycznych, którego dowód znajduje się
w Twierdzeniu~\ref{thm:przestawianie-18}.
To odwołanie naprzód nie tworzy koła: tamten argument
używa przepływu i braku trajektorii między zamienianymi
punktami, a nie homologii Morse'a.
W konstrukcji zmieniamy funkcję, zachowując pole $X$:
nierówność $X\widetilde f<0$ pozostaje prawdziwa
poza punktami krytycznymi, a blisko nich dodajemy
do $f$ tylko stałe.

Jeśli dwa sąsiednie punkty w kolejności wartości
mają indeksy malejące, od wyższego punktu
(o mniejszym indeksie) do niższego nie biegnie
trajektoria przez Stwierdzenie~\ref{prop:wymiar-ms-17}.
Można więc zamienić ich wartości. Skończenie wiele
takich zamian daje funkcję $F$, której wszystkie
punkty indeksu $k$ leżą poniżej wszystkich punktów
indeksu $k+1$. Równe wartości można najpierw rozdzielić
małymi zmianami stałymi blisko punktów krytycznych.
Pole $X$ i jego trajektorie oraz ich znaki pozostają te same.
Jeśli chcemy ponownie używać języka gradientu,
poza zerami kładziemy
\[
 \widetilde g(X,X)=-\dd F(X)>0,\qquad
 \widetilde g(X,\ker\dd F)=0,
\]
a na $\ker\dd F$ wybieramy dowolną dodatnią metrykę.
Blisko zer pozostawiamy $g$; sklejenie metryk podziałem
jedności zachowuje liniową tożsamość
$\widetilde g(X,\cdot)=-\dd F$.
Zatem $X=-\nabla_{\widetilde g}F$.

Wybierzmy regularne poziomy $a_k$ po wszystkich
punktach indeksu $k$, a przed punktami wyższych indeksów.
Niech $N_k=F^{-1}(( -\infty,a_k])$,
$N_{-1}=\varnothing$, $N_n=M$.
Twierdzenie o uchwycie i wycięcie dają
\[
 H_j(N_k,N_{k-1};\Z)=
 \begin{cases}\Z^{c_k},&j=k,\\0,&j\ne k.\end{cases}
\]
Przy pustej grupie indeksów wybieramy po prostu
dwa poziomy w tym samym regularnym paśmie.
Filtracja jest równoważna filtracji szkieletami
modelu CW: na każdym etapie przybywają wyłącznie
komórki wymiaru $k$. Jej kompleks względny ma więc
grupy $C_k=H_k(N_k,N_{k-1})$ i brzeg
\[
 C_k\xrightarrow{\delta}H_{k-1}(N_{k-1})
       \longrightarrow H_{k-1}(N_{k-1},N_{k-2})=C_{k-1}.
\]
Jak w dowodzie twierdzenia o homologii komórkowej,
dokładne ciągi par identyfikują jego homologię z $H_*(M)$.

\emph{Krok 2: konkretni reprezentanci generatorów.}
Mały dysk w $W^u(p)$, zorientowany wybraną orientacją,
przedłużamy przepływem aż do poziomu $a_{k-1}$.
Każdy punkt jego sfery wyjściowej dociera tam:
w przeciwnym razie granicą byłby punkt o indeksie
co najmniej $k$, sprzecznie ze ścisłym spadkiem indeksu.
Gładki czas trafienia jest ograniczony na zwartej sferze.
Otrzymujemy zanurzony dysk $D_p\subset N_k$
z brzegiem $S_p\subset\partial N_{k-1}$.
Reprezentuje generator odpowiadający uchwytowi $p$:
lokalnie jego przestrzeń styczna jest ujemną przestrzenią
hesjanu, a więc jego zorientowany dysk jest rdzeniem
generatora lokalnej homologii względnej.
Przedłużenie przez regularny przepływ dodaje tylko kołnierz.

Dla $q$ indeksu $k-1$ analogicznie przedłużamy mały
dysk stabilny wstecz do $\partial N_{k-1}$.
Dostajemy \emph{współrdzeń} $E_q$ uchwytu $q$
i jego brzeg $B_q\subset\partial N_{k-1}$,
nazywany sferą pasa. Nie może on zatrzymać się
przy innym punkcie indeksu $k-1$, bo nie ma
trajektorii między równymi indeksami.
Normalną orientację $E_q$ wyznacza $W^u(q)$.
W małym modelu uchwytu rdzeń i współrdzeń przecinają
się raz, ze znakiem $+1$ przy zgodnych orientacjach.
Jest to lokalna dualność, która pozwala odczytać
współczynnik generatora w $H_{k-1}(N_{k-1},N_{k-2})$.

\emph{Krok 3: porównanie współczynników.}
Obraz $[D_p]$ przez $\delta$ to $[S_p]$.
Jego współczynnik przy $q$ liczymy jako podpisaną
liczbę przecięć $S_p$ z $B_q$ w $\partial N_{k-1}$.
Dlaczego ten rachunek daje współczynnik względny?
W otoczeniu tubularnym współrdzenia rzut na normalny
dysk, po zgnieceniu jego brzegu, mierzy lokalny stopień.
Wycięcie redukuje go do pary uchwytu i jego części
przyczepianej, gdzie rdzeń daje stopień $1$,
a pozostałe generatory dają $0$.
Jest to dokładnie funkcjonał dualny do generatora $q$.
Poprzeczność pozwala obliczać stopień jako sumę
lokalnych znaków, również bez orientowania całego $M$;
wystarczają orientacja $S_p$ i koorientacja $B_q$.

Punkty $S_p\cap B_q$ odpowiadają bijektywnie
trajektoriom $p\to q$, gdyż każda przecina ten
regularny poziom dokładnie raz. Orientacja brzegu
$D_p$ jest określona przez zewnętrzny kierunek $X$,
więc zgadza się z przyjętą orientacją $S_p$.
Koorientacja $B_q$ jest tą samą, której użyliśmy
w definicji $\epsilon(\gamma)$. Współczynnik brzegu
jest zatem dokładnie sumą w~\eqref{eq:brzeg-morse-17}.
Kompleks filtracji jest izomorficzny z kompleksem Morse'a
przez przypisanie $p\mapsto[D_p]$.
Każdy wybór $f,g$ oblicza w ten sposób homologię $M$,
a odwrócenie orientacji $W^u(p)$ zmienia jedynie znak
odpowiedniego generatora.
\end{proof}

\begin{przyklad}[Kompleks na sferze i torusie]
Na $S^2$ z wysokością
z Przykładu~\ref{prz:morse-sfera-17}
nie ma punktów indeksu $1$,
więc
$CM_2=\Z$, $CM_1=0$,
$CM_0=\Z$, a $\partial=0$.
Homologie to $\Z,0,\Z$.

Na torusie z
Przykładu~\ref{prz:morse-torus-17}
mamy
$CM_2=\Z$, $CM_1=\Z^2$,
$CM_0=\Z$. Znamy
$H_0(T^2;\Z)=\Z$,
$H_1(T^2;\Z)=\Z^2$,
$H_2(T^2;\Z)=\Z$;
po przejściu do $\Q$
liczby punktów krytycznych
są już równe liczbom
Betti'ego. Wzór
\eqref{eq:ranki-morse-17}
zmusza wszystkie
$r_k$ do zera.
Macierze całkowite
o zerowym rzędzie
nad $\Q$ są zerowe,
więc $\partial=0$
również nad $\Z$.
Geometrycznie podpisane
liczby trajektorii między
sąsiednimi indeksami
znoszą się w sumie.
\end{przyklad}

\section{Jak to prowadzi do Smale'a i chirurgii?}
\label{sec:morse-most-17}

Funkcja Morse'a opisuje rozmaitość
przez uchwyty. Dla kobordyzmu
$W$ między $M_-$ i $M_+$
analogicznie wybiera się funkcję
stałą na składowych brzegu;
jej punkty krytyczne dają
uchwyty dołączane do
$M_-\times[0,1]$.
Skoro dołączenie uchwytu
indeksu $\lambda$ zmienia
topologię określonego
obszaru, można pytać,
kiedy sąsiednie uchwyty
się znoszą, kiedy da
się przestawiać ich
kolejność i jakie
przeszkody pozostają.
To są kolejne kroki
na drodze do twierdzenia
o $h$-kobordyzmie
Smale'a i teorii chirurgii.

Na rozmaitości niezwartej
właściwa funkcja Morse'a
z Twierdzenia~\ref{thm:wlasciwa-morse-15}
daje analogiczny opis
\emph{po jednym zwartym
paśmie naraz}. Dla
rozmaitości z brzegiem
trzeba kontrolować
kierunek gradientu
przy brzegu i ewentualne
punkty krytyczne
funkcji ograniczonej
do brzegu. Wyniki
tego rozdziału zostały
sformułowane dla
zamkniętej $M$ właśnie
po to, aby takich
warunków nie ukrywać.

\par\smallskip
{\footnotesize\noindent\emph{Dalsza lektura:}
\href{https://math.stanford.edu/~ralph/morsecourse/biglectures.pdf}{R.\ Cohen, \emph{Morse Theory}, lemat Morse'a, uchwyty i nierówności},
\href{https://people.maths.ox.ac.uk/ritter/morse-cambridge/full.pdf}{A.\ Ritter, \emph{Morse Homology}, trajektorie i ich zliczanie},
\href{https://arxiv.org/abs/math/0411465}{J.\ Weber, \emph{The Morse--Witten complex via dynamical systems}, zwartość, sklejanie i orientacje},
\href{https://ocw.mit.edu/courses/18-965-geometry-of-manifolds-fall-2004/pages/lecture-notes/}{MIT OCW, wykłady z geometrii rozmaitości}.\par}
